\subsubsection*{3.1.5 Final Design Class Diagram}
\addcontentsline{toc}{subsubsection}{3.1.5 Final Design Class Diagram}
\pdfbookmark[2]{3.1.5 Final Design Class Diagram}{p7h69}
Author: 722I0004 \ensuremath{\cdot} AI-assisted (Claude)

The final design class diagram refines the layer views of 3.1.2 with what the sequence diagrams in 3.1.4 and the implementation settled: UC-01 is served by two controllers, so the final diagram is drawn in two parts, one per controller; classes shared by both parts appear in each.

\begin{itemize}
\item \textbf{One diagram, four layers.} View (Flask blueprint functions with their routes), controller, domain and data access are drawn together, so every message in the sequence diagrams can be followed from the screen to the table.

\item \textbf{Only the members this use case calls.} The sequence diagrams fixed which operations each class must offer; the complete member lists stay in the 3.1.1 tables and 3.1.2 views.

\item \textbf{Signatures from the code.} Parameter names and types are read from the source at the release commit by \texttt{models/\allowbreak{}shared/\allowbreak{}final\_\allowbreak{}dcd.\allowbreak{}py}, so the diagram and the program cannot differ.

\item \textbf{Navigability and persistence.} Dashed arrows show which layer calls which; each DAO has a \guillemotleft{}persists\guillemotright{} link to the entity it stores; domain associations keep their multiplicities from the domain model (2.1.7).

\end{itemize}
\begin{landscape}
\begin{figure}[p]
\centering
\includegraphics[width=230mm,height=150mm,keepaspectratio]{\detokenize{models/UC-01/dcd_uc01_final_overview_PartA.png}}
\caption{Figure 3.1.5A. UC-01 final design class diagram --- part A: PlaceOrderController --- Author: 722I0004 \ensuremath{\cdot} AI-assisted (Claude). Vector SVG beside the source (editable source: models/UC-01/dcd\_uc01\_final\_overview\_PartA.dot)}\label{P7-0613}
\end{figure}
\end{landscape}
\begin{landscape}
\begin{figure}[p]
\centering
\includegraphics[width=230mm,height=150mm,keepaspectratio]{\detokenize{models/UC-01/dcd_uc01_final_overview_PartB.png}}
\caption{Figure 3.1.5B. UC-01 final design class diagram --- part B: OrderTrackingController --- Author: 722I0004 \ensuremath{\cdot} AI-assisted (Claude). Vector SVG beside the source (editable source: models/UC-01/dcd\_uc01\_final\_overview\_PartB.dot)}\label{P7-0614}
\end{figure}
\end{landscape}
Behaviour refined between the first design and the final design (each refinement is covered by the tests in 5.1). Every state-changing operation runs in one \texttt{BEGIN IMMEDIATE} transaction and checks the row count of each guarded UPDATE, so a stale page, a guessed ID or a competing user changes nothing.

\begin{longtable}{>{\raggedright\arraybackslash}p{0.4600\linewidth}>{\raggedright\arraybackslash}p{0.4600\linewidth}}
System operation & Refinement in the final design \\
\hline\endfirsthead
System operation & Refinement in the final design \\
\hline\endhead
\texttt{requestQuote} & The delivery radius is checked on the raw distance, not a rounded one (8.004 km is refused). \\
\texttt{requestQuote, confirmCheckout} & Consent compares a fingerprint of the whole quote (hub, address, items, prices, voucher, fees, payment), not only the total; any change asks for a fresh quote. \\
\texttt{confirmCheckout} & Capacity is re-counted inside the write transaction, so two customers cannot both take the last slot. \\
\texttt{viewReceipt} & The receipt says an amount is "owed", never "refunded", because payout is out of scope. \\
\texttt{cancelOrder} & Cancelling an online order records the full captured amount as refund owed. \\
\end{longtable}
Editable sources: \texttt{models/\allowbreak{}UC-01/\allowbreak{}dcd\_\allowbreak{}uc01\_\allowbreak{}final\_\allowbreak{}overview\_\allowbreak{}PartA.\allowbreak{}dot} and \texttt{dcd\_\allowbreak{}uc01\_\allowbreak{}final\_\allowbreak{}overview\_\allowbreak{}PartB.\allowbreak{}dot} (Graphviz), with SVG exports beside them.
